\bookchapter{第二十二章：AI點樣喺手機中工作}

\chapterintro{歡迎嚟到手機最酷嘅“魔術現場”，主角係：人工智能。你每日喺度用嘅“一拍即得”嘅靚相、張口就嚟嘅語音助手、越用越明你嘅輸入法，背後都企住呢位隱形大腦。}

\booksection{一、AI相機：你嘅“後期暗房”搬入咗手機}

而家嘅AI計算攝影，更似係將一個專業修圖師+攝影師團隊直接塞入咗快門鍵入面。喺你撳快門前，AI就已經開始工作啦：掃描場景，識別出“呢個係藍天”、“呢個係人臉”、“呢個係夜景”。

\booksection{二、語音助手：你嘅“聲音祕書”點樣聽明你？}

\begin{enumerate}
  \item \textbf{聽清（語音識別）}：將聲波轉變成文字。
  \item \textbf{聽明（自然語言理解）}：理解意圖，提取關鍵實體。
  \item \textbf{執行與回答}：調用對應App或服務，再通過語音合成"講"畀你聽。
\end{enumerate}

\booksection{三、端側AI vs. 雲端AI}

端側AI將計算留喺手機入面，速度快、隱私高、可離線。雲端AI需要聯網，但能處理更複雜嘅任務。

\booksection{四、AI大模型與手機：你嘅隨身智能助手}

2025-2026年，AI大模型（如DeepSeek R1、ChatGPT-5、文心一言5.0、通義千問3等）嘅爆發深刻改變咗手機嘅使用方式。手機唔再淨係執行指令嘅工具，而係能理解、推理同埋創作嘅智能夥伴。端側AI晶片（NPU算力突破80 TOPS）等大模型推理首次喺手機本地跑通，無需聯網都能用AI。

\booksubsection{乜嘢係AI大模型？}

AI大模型（Large Language Model，LLM）係通過海量文本數據訓練出嚟嘅AI系統，能理解同埋生成自然語言。你可以將佢想像成一個讀過互聯網上幾乎所有文字嘅“超級學霸”，能回答問題、寫文章、翻譯、編程、分析數據。

\booksubsection{手機上嘅AI大模型應用：}

\begin{itemize}
  \item \textbf{對話助手}：DeepSeek、ChatGPT、豆包、Kimi等App，可以隨時提問、翻譯、寫文案、總結文章。比傳統搜索引擎更智能——佢唔淨係畀你連結，而係直接畀你答案。
  \item \textbf{寫作與創作}：等AI幫你寫郵件、總結會議紀要、潤色文章、生成創意文案。
  \item \textbf{圖像生成與編輯}：AI一鍵消除相片中嘅路人、擴展圖片邊界、風格轉換、甚至根據文字描述生成圖片。
  \item \textbf{實時翻譯}：通話實時翻譯、面對面對話翻譯、影相翻譯，出國旅行唔再有語言障礙。
  \item \textbf{學習輔導}：AI可以解釋概念、出題練習、批改功課，係隨身家教。
\end{itemize}

\booksubsection{端側大模型：手機自己跑AI}

最新嘅旗艦手機已經能喺本地運行小型AI大模型（通常7B參數左右），帶嚟嘅變化係：

\begin{itemize}
  \item \textbf{唔需要網絡}：斷網都能用AI功能，飛機上、地鐵入面都得。
  \item \textbf{隱私更安全}：你嘅對話同埋數據唔會上傳到雲端，完全喺手機本地處理。
  \item \textbf{響應更快}：冇網絡延遲，AI回覆幾乎係即時嘅。
  \item \textbf{當前限制}：端側模型嘅能力仲比唔上雲端大模型，複雜任務仍需聯網。但隨住晶片AI算力嘅提升，端側模型會越來越強。
\end{itemize}

\booksubsection{系統級AI集成：}

\begin{itemize}
  \item \textbf{蘋果Apple Intelligence}：iOS 18引入嘅AI系統，集成寫作工具、圖像魔法棒、Siri增強、通知摘要等功能。端側處理優先，複雜任務先上雲端。
  \item \textbf{華為盤古大模型}：集成到鴻蒙系統中，小藝助手獲得大模型能力，支援檔案摘要、智慧識屏、AI修圖等。
  \item \textbf{小米澎湃AI}：系統級AI能力，支援AI搜圖、AI摘要、AI實時字幕、端側大模型對話等。
  \item \textbf{三星Galaxy AI}：通話實時翻譯、筆記AI摘要、圖片生成編輯、搜索圈選等功能。
\end{itemize}

\booksection{五、幾個典型應用場景與實操}

為咗等你更直觀咁感受手機AI嘅“超能力”，下面我哋用五個具體場景，展示AI係點樣喺日常生活中“出手相救”嘅：

\booksubsection{場景 1：實時翻譯——出國旅行嘅隨身翻譯官}

\textbf{實際表現}：喺巴黎點餐唔識法語？iPhone嘅“翻譯”App支援面對面雙人對話模式，雙方各睇自己語言嘅字幕，甚至可以同聲傳譯。安卓陣營嘅三星Galaxy AI通話實時翻譯能等你喺電話入面直接講中文，對方聽到嘅係英文，反之亦然。

\textbf{背後原理}：語音識別（ASR） 機器翻譯（NMT） 語音合成（TTS）三步流水線，端到端延遲已可控制喺1秒內。

\booksubsection{場景 2：AI消除——拯救被路人“毀掉”嘅遊客照}

\textbf{實際表現}：喺景點影咗一張絕美嘅自拍，但背景入面擠滿咗其他遊客？用“魔法消除”圈選路人，手機會自動識別並將路人“塗掉”，再智能填充背景，效果成日令人驚艷。

\textbf{背後原理}：圖像分割（Inpainting）模型先識別並刪除選中區域，然後根據周圍像素同埋語義理解“腦補”出合理嘅填充內容。呢背後嘅AI模型訓練數據量高達數百萬張圖片。

\booksubsection{場景 3：AI修圖——一鍵變“攝影大師”}

\textbf{實際表現}：影咗一張陰天嘅灰濛濛風景相？AI“魔法天空”可以一鍵將天空換成晴空、晚霞甚至星空。AI擴圖可以等豎構圖變成橫構圖，手機“腦補”出畫面外嘅場景。

\textbf{背後原理}：基於生成式AI（Generative AI）嘅擴散模型（Diffusion Model），理解圖像嘅語義內容後生成新像素。呢個係當前AI圖像處理嘅“最前沿”。

\booksubsection{場景 4：檔案處理——辦公族嘅效率倍增器}

\textbf{實際表現}：會議結束後，將會議錄音掟畀AI，佢能自動整理出結構化嘅會議紀要、提取待辦事項、生成摘要。一份10頁嘅英文PDF報告，AI可以喺10秒內畀你提煉出核心結論。

\textbf{背後原理}：長上下文大模型（Long Context LLM）可以“讀完”數萬字嘅檔案，進行多輪推理同埋資訊抽取。Apple Intelligence、華為小藝、WPS AI、騰訊元寶等都具備此能力。

\booksubsection{場景 5：AI搜圖與識物——你嘅“視覺搜索引擎”}

\textbf{實際表現}：喺街度睇到一種唔識嘅花，相機一開，AI即刻話你知“呢個係木槿花”。相簿入面翻搵舊年喺海邊影嘅相，搜索“大海”，AI自動將所有含海嘅相搵出嚟，唔使手動翻。

\textbf{背後原理}：計算機視覺（CV）模型對圖片內容進行向量化（Embedding），再同文字查詢進行相似度匹配。呢種“多模態搜索”係手機相簿管理嘅一大飛躍。

\booksection{六、提示詞優化：等 AI 聽你指揮嘅“溝通公式”}

AI 嘅回答質量，好大程度上取決於你點樣“提問”。同一個問題，問法唔同，AI 畀嘅答案天差地別。下面我哋介紹一套\textbf{最實用、最容易上手嘅“角色-背景-任務-限制”溝通思路}，並配以大量實例。

\begin{tcbinfo}{點解需要框架？}
好多人用 AI 似用百度噉——敲幾個關鍵詞入去。但 AI 唔係搜索引擎，佢係一個“閲讀理解後回答”嘅語言模型。\textbf{畀佢更多上下文、講清你要乜、限制佢唔好跑偏}，佢就能畀你精準嘅結果。框架就係將你腦入面嘅“好答案長點樣”翻譯成 AI 能理解嘅指令。

\end{tcbinfo}

\booksubsection{角色（Role）：話畀 AI 聽佢係邊個}

畀 AI 一個具體嘅“身份”，佢就會調用對應領域嘅專業知識同埋表達方式。\textbf{角色設定得越具體，輸出越專業}。

\booksubsection{實用示例：}

\begin{itemize}
  \item "你係一位\textbf{有 10 年經驗嘅小米手機維修工程師}。"（技術問題）
  \item "你係一位\textbf{擅長將複雜概念講簡單嘅中學物理老師}。"（科普問題）
  \item "你係一位\textbf{審美好、讀得多、有自己風格嘅圖書編輯}。"（寫作問題）
  \item "你係一位\textbf{每日畀上班族做早餐嘅烹飪博主}。"（生活問題）
  \item "你係一位\textbf{耐心、溫同、永遠唔嘲笑新手嘅中醫}。"（健康問題）
\end{itemize}

\textbf{點解有效？}角色設定咗 AI 嘅\textbf{知識範圍}、\textbf{表達風格}同埋\textbf{專業深度}。冇角色嘅 AI 係個“全才但乜都唔精”嘅通才；有咗角色，佢能“進入”嗰個身份，似嗰位專業人士噉思考。

\booksubsection{背景（Background）：話畀 AI 聽你嘅情況}

將“你係邊個、你嘅現狀、你嘅約束條件”提前講清楚，AI 就能畀出\textbf{真正適合你}嘅答案，而唔係泛泛而談。

\booksubsection{實用示例：}

\begin{itemize}
  \item "我嘅手機係\textbf{小米 14，買咗半年，最近打機時發燙嚴重}。"（設備情況）
  \item "我係一個 60 歲嘅退休阿姨，眼唔太好，唔太識用智能手機。"（用戶畫像）
  \item "我正在\textbf{準備 2026 年 6 月嘅英語六級考試，基礎一般}。"（學習背景）
  \item "我開咗一間\textbf{三線城市嘅小咖啡店，月流水 3 萬，請咗兩個員工}。"（商業背景）
  \item "我\textbf{啱啱生完BB喺屋企坐月，睡眠嚴重不足}。"（生活階段）
\end{itemize}

\textbf{點解有效？}背景資訊等 AI 知道\textbf{“為邊個服務”}同埋\textbf{“喺乜嘢條件下”}。同一道題，畀學生同畀家長講法唔同；畀城入麪人同畀鄉鎮人建議都唔同。背景缺失，AI 淨係可以畀“最大公約數”嘅答案，往往唔痛唔癢。

\booksubsection{任務（Task）：明確話畀 AI 聽做乜}

呢一步要用\textbf{動詞開頭}（寫、列、對比、翻譯、解釋、總結……），將任務描述得\textbf{具體、可執行}。含糊嘅任務淨係會換嚟含糊嘅答案。

\booksubsection{實用示例：}

\begin{itemize}
  \item "請\textbf{列出 5 個我應該即刻關閉嘅權限}，並說明原因。"（動作明確）
  \item "請\textbf{將下面呢段英文郵件翻譯成中文}，保持商務口吻。"（任務具體）
  \item "請\textbf{對比呢四款 3000 元價位嘅手機}，畀出一張優缺點對比表。"（格式明確）
  \item "請\textbf{用小學生能明嘅話解釋'乜嘢係 5G'}。"（風格明確）
  \item "請\textbf{畀我嘅早餐店寫三條短影片文案}，分別針對上班族、情侶、帶仔女嘅家長。"（細分明確）
\end{itemize}

\textbf{點解有效？}任務部分係 AI 行動嘅“指令中樞”。\textbf{動詞越清晰、要求越具體，AI 越知道“交付乜嘢”}。模糊嘅“搞下”、“整個”、“睇下”係 AI 大忌——佢會揣測你嘅意思，結果成日跑偏。

\booksubsection{限制（Constraints）：話畀 AI 聽唔好做乜}

“限制”聽起嚟似“約束”，其實係\textbf{等 AI 收斂到最適合你嘅輸出}嘅利器。常見維度有：字數、格式、語氣、避免項、必含項。

\booksubsection{實用示例：}

\begin{itemize}
  \item "\textbf{唔超過 200 字，分點列出}。"（長度+格式）
  \item "\textbf{語氣要親切自然，似朋友傾偈}，唔好用'首先、其次、再次'呢種官話。"（語氣）
  \item "\textbf{唔好使用專業術語}，遇到專業名詞要順便解釋。"（用詞限制）
  \item "\textbf{每段唔超過 3 行}，方便手機閲讀。"（排版）
  \item "\textbf{必須包含：使用場景、優缺點、我嘅預算匹配度}。"（必含項）
  \item "\textbf{唔好畀具體投資建議}，淨係解釋原理。"（邊界）
\end{itemize}

\textbf{點解有效？}冇限制，AI 容易“自由發揮”——要麼太學術、要麼太長、要麼跑題。明確嘅限制=明確嘅“交付標準”，雙方對齊預期，結果先至可控。

\booksubsection{提示詞質量評估：好提示詞嘅 5 個特徵}

寫好提示詞唔係一次性嘅事，\textbf{寫完先自查呢五條}，再發畀 AI：

\begin{enumerate}
  \item \textbf{具體而非抽象}：唔係"寫篇文章"，而係"寫一篇 800 字、3 個小節、面向新手媽媽嘅手機選購指南"。
  \item \textbf{包含背景而非孤立}：話畀 AI 聽"我係邊個、點解問呢個"，佢先至能畀出"為你定製"嘅答案。
  \item \textbf{有明確輸出格式}：係表格、列表、段落、定係 JSON？提前講。
  \item \textbf{限制到位}：字數、風格、避項，三選一就有明顯提升。
  \item \textbf{可被 AI 驗證}：問完之後，等 AI 複述一遍你嘅要求；講唔清楚就係冇講清楚。
\end{enumerate}

\booksubsection{常見問題與解法}

\begin{center}
\begin{tabularx}{\linewidth}{|X|X|}
\hline
\textbf{症狀} & \textbf{解法} \tabularnewline
\hline
回答太泛、唔實用 & 加背景（你係邊個/乜嘢場景） \tabularnewline
回答太學術/太官方 & 加語氣限制 + 拒絕"官話"清單 \tabularnewline
回答跑題/答非所問 & 將任務用動詞重新講一遍，並要求"淨係回答呢個" \tabularnewline
回答太長睇唔完 & 硬性限制字數（200 字內） \tabularnewline
AI 一本正經胡說 & 加"如不確定請標註'我不確定'" + 重要事實自查 \tabularnewline
第一版唔滿意 & 唔好重寫，喺原提示詞後追加"再…一啲" \tabularnewline
\hline
\end{tabularx}
\end{center}

\booksubsection{三個完整場景案例}

\begin{tcbinfo}{場景一：生活助手（早餐搭配建議）}
\textbf{提示詞}：

你係一位擅長營養搭配嘅早餐博主（\textbf{角色}）。我係一位上班族媽媽，每日朝早淨係得 15 分鐘做早餐，屋企有 6 歲嘅細路同埋 60 歲嘅奶奶（\textbf{背景}）。請\textbf{列出 7 日唔重樣嘅早餐搭配}（\textbf{任務}）。每份早餐要求：(1) 準備時間唔超過 15 分鐘 (2) 營養均衡（蛋白+碳水+蔬果）(3) 老少皆宜 (4) 食材喺普通超市能買到（\textbf{限制}）。

\textbf{點解有效}：角色定專業度，背景畀三類人共性，任務明確交付物，4 條限制確保可執行。

\end{tcbinfo}

\begin{tcbinfo}{場景二：學習輔助（英語六級備考）}
\textbf{提示詞}：

你係一位\textbf{經驗豐富嘅英語六級老師}（角色）。我係一名普通本科生，\textbf{四級 480 分，六級目標 500，仲有 3 個月}（背景）。請\textbf{為我制定一份可執行嘅 12 週備考計劃}（任務）。要求：(1) 每週明確聽說讀寫嘅時間分配 (2) 畀出每個月嘅重點突破方向 (3) 推薦 3 本合適嘅複習資料（中文說明理由） (4) 表格形式輸出（限制）。

\textbf{點解有效}：角色定專業背景，背景畀出分數同埋時間錨點，任務明確“分階段計劃”，限制要求“表格化”同埋“可執行”。

\end{tcbinfo}

\begin{tcbinfo}{場景三：工作效率（會議紀要整理）}
\textbf{提示詞}：

你係一位\textbf{資深行政助理}（角色）。我啱啱開完一個 1 小時嘅部門會議，\textbf{討論 Q2 營銷推廣方案}（背景）。\textbf{請將以下會議記錄整理成結構化紀要}（任務）。要求：(1) 分「決議事項 / 待辦事項 / 風險與異議」三部分 (2) 待辦事項標出負責人同截止時間（如未提到則寫「待確認」） (3) 總字數控制喺 600 字內 (4) 用 markdown 格式輸出（限制）。

\textbf{點解有效}：角色定行政視角，背景畀會議主題，任務明確結構化輸出，4 條限制等結果可以直接重用。

\end{tcbinfo}

\textbf{進階玩法：等 AI 成為你嘅「私人助理」}\par

試嚇呢啲進階用法，體驗 AI 嘅真正實力：

\begin{itemize}
  \item \textbf{多輪對話優化}：對 AI 嘅回答唔滿意時，可以講「呢個太專業喇，請用小學生聽得明嘅話再解釋一次」。
  \item \textbf{鏈式思考（CoT）}：等 AI「一步步思考」，回答質素通常更高：「請一步步分析呢道數學題」。
  \item \textbf{多 AI 協作}：用 A 寫初稿，B 潤色，C 翻譯——組合使用效果更佳。
  \item \textbf{系統提示詞}：喺 AI 設定中預設你嘅身份、偏好、忌諱等，等回答更貼合你嘅需求。
\end{itemize}

\booksection{七、AI嘅侷限與挑戰：保持清醒嘅判斷力}

AI雖然強大，但絕非「無所不能」。清醒地認識佢嘅侷限，先至可以更好咁使用佢：

\begin{itemize}
  \item \textbf{「幻覺」問題}：AI可能一本正經咁胡說八道（專業術語叫Hallucination）。佢會編造唔存在嘅歷史事件、錯誤嘅引言、虛假嘅統計數據。重要資訊必須人工核實。
  \item \textbf{知識截止}：大多數AI嘅訓練數據有截止日期（如2024年初），佢唔知之後發生嘅事。
  \item \textbf{推理能力有限}：複雜嘅多步推理、數學證明、邏輯陷阱，AI仍然容易出錯。
  \item \textbf{私隱風險}：喺雲端AI服務中，你嘅對話可能被用於模型訓練。敏感資訊應優先使用端側AI。
  \item \textbf{偏見與價值觀}：AI可能帶有訓練數據中嘅偏見，需要保持獨立思考。
\end{itemize}

\begin{tcbinfo}{AI使用小貼士}
\begin{itemize}
  \item AI生成嘅內容可能唔準確（「幻覺」問題），重要資訊務必核實。
  \item 唔好向AI輸入個人私隱資訊（身份證號、密碼、銀行資訊等）。
  \item 善用提示詞（Prompt）：越具體、越清晰嘅描述，AI嘅回答質素越高。
  \item AI係工具唔係權威，保持獨立思考同判斷力。
  \item 遇到關鍵決策（醫療、法律、財務），諮詢專業人士而唔係依賴AI。
\end{itemize}

\end{tcbinfo}

\begin{tcbpractice}{嚟挑戰：三道AI實戰題}
\begin{enumerate}
  \item \textbf{挑戰一：複雜指令}：對語音助手講：「明早7點提醒我帶遮，同埋將今日下晝3點嘅日曆改到4點。」睇嚇佢能否一次完成多步操作。
  \item \textbf{挑戰二：AI修圖}：影一張有路人嘅遊客照，使用「魔法消除」或「AI橡皮擦」功能，將路人清除乾淨。觀察AI對背景嘅「腦補」效果。
  \item \textbf{挑戰三：提示詞優化}：先用簡單提示詞「寫一段旅遊攻略」問AI，再用CRISPE框架寫一個詳細提示詞。比較兩次回答嘅質素差異。
\end{enumerate}

\end{tcbpractice}

\begin{tcbinfo}{案例 2：換機數據遷移完整流程}
\begin{itemize}
  \item \textbf{場景}：iPhone 13 小米 15 跨品牌換機。
  \item \textbf{步驟}：(1) 舊機 iCloud/微信自帶備份 (2) 退出 Apple ID + 關閉「查找」 (3) 新機用「小米換機」App 掃碼 (4) 選擇遷移項 (5) 微信單獨用自帶「聊天記錄遷移」 (6) 驗證關鍵數據（通訊錄/相冊/微信） (7) 舊機加密擦除。
  \item \textbf{預期結果}：數據零丟失，舊機可安全轉賣。
\end{itemize}

\end{tcbinfo}
